\chapter{Bentuk Diferensial dan Teorema Stokes}\label{ch:b3:forms}

Satu teorema analisis, selama dua abad, telah menyerap semua
teorema sejenisnya: teorema dasar kalkulus,
Green--Riemann (yang dibuktikan pada jilid Tahun~2), teorema
divergensi Gauss, teorema curl Kelvin--Stokes --- masing-masingnya
mengatakan bahwa integral suatu turunan atas sebuah daerah sama dengan
integral objek semulanya atas batasnya. Bahasa
\emph{\omterm{def:b3:forms:diffform}{bentuk diferensial}} menjadikan semuanya satu pernyataan,
$\int_M\dd\omega = \int_{\partial M}\omega$, dan menjadikan pernyataan
itu terbuktikan dalam satu tarikan. Bab ini membangun
bahasanya dengan jujur --- aljabar multilinear berselang-seling,
\omterm{def:b3:forms:d}{turunan eksterior}, \omterm{def:b3:forms:pullback}{tarikan balik}, \omterm{def:b3:forms:orientation}{orientasi}, pengintegralan pada
submanifold \cref{ch:b3:submanifolds} --- membuktikan teorema
Stokes, lalu mencairkan cek pertamanya: teorema integral
klasiknya, \omterm{def:b3:forms:winding}{bilangan lilitan} yang diam-diam menjalankan
\cref{ch:b3:residues}, dan, pada soal akhir pekannya, teorema
titik tetap Brouwer. Sepanjang bab ini, \emph{mulus} berarti $\mathcal
C^\infty$; setiap pemetaan dan bentuk bersifat mulus kecuali
dinyatakan lain. Ini tak mengorbankan keumuman yang layak dimiliki pada
aras ini dan justru membebaskan tangan.

\section{Aljabar multilinear berselang-seling}

\begin{definition}\label{def:b3:forms:alternating}
Misalkan $E$ ruang vektor real berdimensi $n$. Sebuah
\emph{bentuk berselang-seling $k$-linear}\index{bentuk berselang-seling}
pada $E$ adalah pemetaan $\alpha\colon E^k \to \R$ yang linear pada setiap
variabelnya, dengan $\alpha(v_1, \dots, v_k) = 0$ setiap kali dua
argumennya sama. Ruangnya ditulis $\Lambda^k E^*$;
dan secara konvensi $\Lambda^0E^* = \R$. Sifat berselang-seling itu memaksa
keantisetangkupan: menukar dua argumen mengubah tandanya
(uraikanlah $\alpha(\dots, v + w, \dots, v + w, \dots) = 0$), dan
lebih umum $\alpha(v_{\sigma(1)}, \dots, v_{\sigma(k)}) =
\varepsilon(\sigma)\,\alpha(v_1, \dots, v_k)$ bagi setiap
permutasi $\sigma$.
\end{definition}

\begin{example}
Pada $E = \R^n$: $\Lambda^1E^* = E^*$ adalah \omterm{def:b3:banach:operator}{ruang dualnya};
determinan pada basis kanoniknya merupakan \omterm{def:b3:forms:alternating}{bentuk berselang-seling}
$n$-linear, dan \cref{prop:b3:forms:basis} akan menunjukkan bahwa ia
merentang $\Lambda^nE^*$ --- yakni alasan mendalam mengapa determinan
tunggal sampai penskalaan. Untuk $k > n$: $\Lambda^kE^* = \{0\}$, sebab
$k$ vektor pasti bergantung, dan menguraikan salah satunya sepanjang yang lain
membunuh $\alpha$ menurut sifat berselang-selingnya.
\end{example}

\begin{definition}\label{def:b3:forms:wedge}
Untuk $\ell_1, \dots, \ell_k \in E^*$, \emph{hasil kali
eksterior}\index{hasil kali eksterior} bagi keluarga itu adalah \omterm{def:b3:forms:alternating}{bentuk
berselang-seling} $k$-linear
\[
(\ell_1 \wedge \dots \wedge \ell_k)(v_1, \dots, v_k)
= \det\bigl(\ell_i(v_j)\bigr)_{1 \leq i, j \leq k} .
\]
Sifat berselang-seling dan multilinearitasnya adalah sifat determinan pada
kolomnya.
\end{definition}

\begin{proposition}[Basis $\Lambda^k$]\label{prop:b3:forms:basis}
Misalkan $(e_1, \dots, e_n)$ basis $E$ dengan basis \omterm{def:b3:banach:operator}{dual}
$(e_1^*, \dots, e_n^*)$. Bentuk
\[
e_I^* = e_{i_1}^* \wedge \dots \wedge e_{i_k}^*,
\qquad I = \{i_1 < \dots < i_k\} \subseteq \{1, \dots, n\},
\]
membentuk basis $\Lambda^kE^*$; sehingga $\dim\Lambda^kE^* =
\binom nk$. Secara eksplisit, $\alpha =
\sum_{\abs I = k}\alpha(e_{i_1}, \dots, e_{i_k})\,e_I^*$.
\end{proposition}

\begin{proof}
\emph{Membangun.} Misalkan $\alpha \in \Lambda^kE^*$ dan $\beta =
\sum_I\alpha(e_I)\,e_I^*$, dengan $\alpha(e_I)$ menyingkat
$\alpha(e_{i_1}, \dots, e_{i_k})$. Kedua ruasnya $k$-linear
dan berselang-seling, sehingga keduanya sama begitu keduanya sama pada semua
$k$-tupel $(e_{j_1}, \dots, e_{j_k})$ dengan $j_1 < \dots <
j_k$ (sebab multilinearitasnya menyusutkan ke tupel vektor basis,
sedangkan sifat berselang-selingnya ke tupel yang naik sejati). Sedangkan $e_I^*(e_{j_1},
\dots, e_{j_k}) = \det(e_{i_r}^*(e_{j_s})) = \delta_{IJ}$:
untuk $I = J$ matriksnya identitas; sedangkan untuk $I \neq J$ ada
baris yang nol. Jadi $\beta(e_J) = \alpha(e_J)$ bagi setiap $J$:
$\beta = \alpha$. \emph{Kebebasan.} Jika $\sum_I c_Ie_I^* = 0$,
maka mengevaluasinya pada $(e_{j_1}, \dots, e_{j_k})$ memberikan $c_J = 0$.
\end{proof}

\begin{definition}\label{def:b3:forms:wedgegeneral}
\omterm{def:b3:forms:wedge}{Hasil kali eksteriornya} diperluas menjadi pemetaan bilinear
$\Lambda^kE^* \times \Lambda^\ell E^* \to
\Lambda^{k+\ell}E^*$, yang ditentukan oleh kebilinearannya dan
$(e_I^*) \wedge (e_J^*) = e_I^* \wedge e_J^*$ (gabungkanlah
lalu urutkan kembali; hasil kalinya $0$ jika $I \cap J \neq
\varnothing$). Ia asosiatif, dan bersifat
\emph{antikomutatif berderajat}:
\[
\beta \wedge \alpha = (-1)^{k\ell}\,\alpha \wedge \beta
\qquad (\alpha \in \Lambda^k, \ \beta \in \Lambda^\ell).
\]
\end{definition}

\begin{proof}[\omnameProof]
Kedua sifatnya diperiksa pada unsur basisnya lalu diperluas lewat
kebilinearannya. Untuk keasosiatifannya: kedua pengurungan
$e_I^* \wedge e_J^* \wedge e_K^*$ sama dengan \omterm{def:b3:forms:wedge}{hasil kali eksterior} keluarga
bentuk-$1$ yang digabungkan, menurut rumus determinan
\cref{def:b3:forms:wedge} (yakni uraian Laplace per blok).
Untuk aturan tandanya: memindahkan masing-masing $\ell$ faktor $\beta$
melewati $k$ faktor $\alpha$ memakan satu tanda per
transposisi bertetangga (yakni penukaran dua baris determinannya), sehingga
$(-1)^{k\ell}$ seluruhnya.
\end{proof}

\begin{proposition}[Tarikan balik, kasus linear]
\label{prop:b3:forms:linearpullback}
Pemetaan linear $u\colon E \to F$ menginduksikan, bagi setiap $k$,
pemetaan linear $u^*\colon \Lambda^kF^* \to \Lambda^kE^*$,
$(u^*\alpha)(v_1, \dots, v_k) = \alpha(u(v_1), \dots,
u(v_k))$. Ia memenuhi $u^*(\alpha \wedge \beta) = u^*\alpha
\wedge u^*\beta$ dan $(u \circ w)^* = w^* \circ u^*$.
Lebih jauh:
\begin{enumerate}
  \item Jika $\dim E = n$ dan $u\colon E \to E$, maka pada
        garis $\Lambda^nE^*$: $u^*\alpha = (\det u)\,\alpha$.
  \item Jika $\operatorname{rk}u < k$, maka $u^* = 0$ pada
        $\Lambda^kF^*$.
\end{enumerate}
\end{proposition}

\begin{proof}
Identitas fungtorialnya langsung dari definisinya
(untuk aturan hasil kalinya, periksalah pada \omterm{def:b3:forms:wedge}{hasil kali eksterior} bentuk-$1$ dengan
rumus determinannya --- $\det(\ell_i(uv_j)) =
\det((u^*\ell_i)(v_j))$ --- lalu perluaslah secara bilinear).
(1) Pemetaan $u^*$ membawa $\Lambda^nE^*$ yang berdimensi satu
(\cref{prop:b3:forms:basis}) ke dirinya sendiri, sehingga $u^*\alpha =
c\,\alpha$ dengan $c$ yang tak bergantung pada $\alpha \neq 0$; lalu mengujinya
pada $\alpha = e_1^* \wedge \dots \wedge e_n^*$ dan $(v_j) =
(e_j)$ memberikan $c = \det(e_i^*(ue_j)) = \det u$.
(2) Untuk $v_1, \dots, v_k \in E$, vektor $u(v_1), \dots,
u(v_k)$ terletak pada peta $u$, yang berdimensi $< k$: sehingga semuanya
bergantung linear, sedangkan \omterm{def:b3:forms:alternating}{bentuk berselang-seling} lenyap pada keluarga
yang bergantung (uraikanlah vektor yang bergantung itu sepanjang yang
lain).
\end{proof}

\section{Bentuk diferensial dan turunan eksterior}

\begin{definition}\label{def:b3:forms:diffform}
Misalkan $U \subseteq \R^n$ terbuka. Sebuah \emph{bentuk
diferensial}\index{bentuk diferensial} berderajat $k$ pada $U$ adalah pemetaan mulus
$\omega\colon U \to \Lambda^k(\R^n)^*$; dan pada basis
\cref{prop:b3:forms:basis} (dengan $\dd x_i$ ditulis bagi
$e_i^*$),
\[
\omega = \sum_{\abs I = k} a_I\,\dd x_I,
\qquad
\dd x_I = \dd x_{i_1} \wedge \dots \wedge \dd x_{i_k},
\]
dengan koefisien mulus $a_I \in \mathcal C^\infty(U)$.
Ruangnya adalah $\Omega^k(U)$; dan $\Omega^0(U) = \mathcal
C^\infty(U)$. Bentuk-$0$ adalah fungsi; bentuk-$1$ adalah medan
bentuk linear (misalnya diferensial $\dd f$ sebuah
fungsi); sedangkan bentuk-$n$ berupa $a\,\dd x_1\wedge\dots\wedge\dd
x_n$, yakni integran alami \cref{ch:b3:product}.
\end{definition}

\begin{definition}[Turunan eksterior]\label{def:b3:forms:d}
Yang disebut \emph{turunan eksterior}\index{turunan eksterior} adalah
pemetaan linear $\dd\colon \Omega^k(U) \to \Omega^{k+1}(U)$
yang didefinisikan oleh
\[
\dd\Bigl(\sum_I a_I\,\dd x_I\Bigr)
= \sum_I \dd a_I \wedge \dd x_I
= \sum_I\sum_{j=1}^n \frac{\partial a_I}{\partial
x_j}\,\dd x_j \wedge \dd x_I .
\]
Pada bentuk-$0$ ia adalah diferensial biasanya.
\end{definition}

\begin{theorem}\label{thm:b3:forms:dsquare}
(a) $\dd(\omega \wedge \eta) = \dd\omega \wedge \eta +
(-1)^k\,\omega \wedge \dd\eta$ bagi $\omega \in \Omega^k$
(yakni \emph{aturan Leibniz berderajat}).
(b) $\dd \circ \dd = 0$.\index{d kuadrat nol@$\dd^2 = 0$}
\end{theorem}

\begin{proof}
(a) Menurut kebilinearannya, cukup ditangani $\omega = a\,\dd
x_I$, $\eta = b\,\dd x_J$. Maka $\omega \wedge \eta =
ab\,\dd x_I \wedge \dd x_J$ dan
\[
\dd(\omega\wedge\eta) = (b\,\dd a + a\,\dd b) \wedge \dd x_I
\wedge \dd x_J
= (\dd a \wedge \dd x_I) \wedge (b\,\dd x_J) +
a\,\dd b \wedge \dd x_I \wedge \dd x_J,
\]
sedangkan memindahkan bentuk-$1$ $\dd b$ melewati $k$ faktor $\dd
x_I$ memakan $(-1)^k$ (\cref{def:b3:forms:wedgegeneral}): sehingga
suku keduanya adalah $(-1)^k\,\omega \wedge \dd\eta$.
(b) Untuk sebuah fungsi: $\dd(\dd a) = \sum_{i,j}
\frac{\partial^2a}{\partial x_j\partial x_i}\,\dd x_j \wedge
\dd x_i$. Koefisiennya setangkup terhadap $(i,j)$ menurut
teorema Schwarz tentang turunan parsial campurannya (yang dibuktikan pada jilid
Tahun~2; sebab $a$ bersifat $\mathcal C^\infty$), sedangkan $\dd x_j \wedge
\dd x_i$ antisetangkup: sehingga dengan memasangkan suku $(i,j)$ dan
$(j,i)$, semuanya meniadakan diri. Untuk $\omega = \sum
a_I\dd x_I$ yang umum: $\dd\dd\omega = \sum\dd(\dd a_I \wedge \dd x_I)
= \sum(\dd\dd a_I\wedge\dd x_I - \dd a_I\wedge\dd(\dd x_I))$
menurut (a), dan kedua sukunya lenyap (sebab $\dd(\dd x_I) = 0$ karena
koefisiennya konstan).
\end{proof}

\begin{definition}[Tarikan balik]\label{def:b3:forms:pullback}
Misalkan $\varphi\colon U \to V$ mulus ($U \subseteq \R^m$,
$V \subseteq \R^n$ terbuka). Yang disebut \emph{tarikan balik}\index{tarikan balik},
$\varphi^*\colon \Omega^k(V) \to \Omega^k(U)$, didefinisikan
secara titik demi titik oleh tarikan balik linear sepanjang diferensialnya:
$(\varphi^*\omega)_x = (D\varphi(x))^*\,\omega_{\varphi(x)}$.
Secara konkret, $\varphi^*$ mensubstitusi: $\varphi^*f = f \circ
\varphi$ pada fungsi, $\varphi^*(\dd y_i) =
\dd\varphi_i = \sum_j\frac{\partial\varphi_i}{\partial
x_j}\dd x_j$, dan $\varphi^*(a\,\dd y_{i_1}\wedge\dots\wedge
\dd y_{i_k}) =
(a\circ\varphi)\,\dd\varphi_{i_1}\wedge\dots\wedge
\dd\varphi_{i_k}$.
\end{definition}

\begin{theorem}\label{thm:b3:forms:pullback}
(a) $\varphi^*(\omega \wedge \eta) = \varphi^*\omega \wedge
\varphi^*\eta$ dan $(\psi \circ \varphi)^* = \varphi^* \circ
\psi^*$.
(b) $\varphi^*(\dd\omega) = \dd(\varphi^*\omega)$: yakni
\omterm{def:b3:forms:d}{turunan eksteriornya} komutatif dengan setiap substitusi mulus
--- yaitu identitas yang menjadikannya turunan \emph{yang sebenarnya}
bagi teori ini.
(c) Jika $\varphi\colon U \to V$ mulus antara \omterm{def:b3:topology:topology}{himpunan terbuka}
$\R^n$ dan $\omega = a\,\dd y_1 \wedge \dots \wedge \dd
y_n$, maka $\varphi^*\omega = (a \circ
\varphi)\,\det\bigl(D\varphi\bigr)\,\dd x_1 \wedge \dots
\wedge \dd x_n$.
\end{theorem}

\begin{proof}
(a) Pernyataan titik demi titik yang menyangkut \omterm{def:b3:forms:pullback}{tarikan balik} linearnya
(\cref{prop:b3:forms:linearpullback}), ditambah aturan rantainya
$D(\psi\circ\varphi)(x) = D\psi(\varphi(x))\,D\varphi(x)$.
(b) Untuk bentuk-$0$ $f$: $\varphi^*(\dd f) = \dd f \circ
D\varphi = \dd(f \circ \varphi)$ adalah aturan rantainya. Untuk
$\omega = a\,\dd y_I$: dengan memakai (a), $\varphi^*\omega =
(a\circ\varphi)\,\dd\varphi_{i_1}\wedge\dots\wedge
\dd\varphi_{i_k}$, sehingga menurut aturan Leibniz
(\cref{thm:b3:forms:dsquare}(a)) dan $\dd\dd\varphi_{i_r} =
0$:
\[
\dd(\varphi^*\omega) = \dd(a\circ\varphi) \wedge
\dd\varphi_{i_1}\wedge\dots\wedge\dd\varphi_{i_k}
= \varphi^*(\dd a) \wedge \varphi^*(\dd y_I)
= \varphi^*(\dd a \wedge \dd y_I) =
\varphi^*(\dd\omega).
\]
(c) Titik demi titik, ini persis $u^*\alpha = (\det u)\alpha$
pada bentuk berderajat tertinggi (\cref{prop:b3:forms:linearpullback}(1)
dengan $u = D\varphi(x)$).
\end{proof}

\begin{example}[Koordinat kutub]\label{ex:b3:forms:polar}
Untuk $\varphi(r, \theta) = (r\cos\theta, r\sin\theta)$:
$\varphi^*\dd x = \cos\theta\,\dd r -
r\sin\theta\,\dd\theta$, $\varphi^*\dd y = \sin\theta\,\dd r
+ r\cos\theta\,\dd\theta$, sehingga
\[
\varphi^*(\dd x \wedge \dd y) =
(\cos\theta\,\dd r - r\sin\theta\,\dd\theta) \wedge
(\sin\theta\,\dd r + r\cos\theta\,\dd\theta)
= r\,\dd r \wedge \dd\theta,
\]
yakni Jacobi \cref{ex:b3:product:polar} yang muncul lewat aljabar
murni --- tanpa teori \omterm{def:b3:measure:measure}{ukuran}. \cref{exo:b3:forms:9} mengubah
catatan ini menjadi sebuah pernyataan: bahwa bagi integral berorientasi,
rumus penggantian variabelnya \emph{adalah} rumus \omterm{def:b3:forms:pullback}{tarikan baliknya}.
\end{example}

\section{Bentuk tertutup dan eksak; lema Poincaré}

\begin{definition}\label{def:b3:forms:closedexact}
Bentuk $\omega \in \Omega^k(U)$ disebut \emph{tertutup}\index{bentuk tertutup}
jika $\dd\omega = 0$, dan \emph{eksak}\index{bentuk eksak} jika
$\omega = \dd\eta$ bagi suatu $\eta \in \Omega^{k-1}(U)$ (yakni sebuah
\emph{primitif} $\omega$). Eksak $\Rightarrow$ tertutup menurut
$\dd^2 = 0$; sedangkan konversnya adalah pertanyaan tentang
\emph{rupa} $U$.
\end{definition}

\begin{example}[Bentuk sudut]\label{ex:b3:forms:angular}
Pada $U = \R^2 \setminus \{0\}$,
\[
\omega_\theta = \frac{x\,\dd y - y\,\dd x}{x^2 + y^2}
\]
bersifat \omterm{def:b3:forms:closedexact}{tertutup} (lewat perhitungan langsung:
\cref{exo:b3:forms:4}) tetapi tak \omterm{def:b3:forms:closedexact}{eksak}: sebab integralnya sepanjang
lingkaran satuannya adalah $2\pi \neq 0$, sedangkan integral \omterm{def:b3:forms:closedexact}{bentuk eksak}
sepanjang kurva \omterm{def:b3:forms:closedexact}{tertutup} lenyap (\cref{prop:b3:forms:exactloop}).
Secara lokal, $\omega_\theta = \dd\theta$ bagi sembarang penentuan
mulus $\theta$ atas sudut kutubnya --- dari sanalah
namanya dan halangannya: bahwa tak ada penentuan semacam itu pada
seluruh $U$. Satu bentuk inilah yang menjalankan \omterm{def:b3:forms:winding}{bilangan lilitan}
(\cref{sec:b3:forms:winding}) dan, lewatnya, teorema \omterm{def:b3:residues:singularities}{residu}
\cref{ch:b3:residues}.
\end{example}

\begin{theorem}[Lema Poincaré]\label{thm:b3:forms:poincare}
Misalkan $U \subseteq \R^n$ terbuka dan \emph{berbentuk bintang}
terhadap $0$. Maka setiap bentuk-$k$ \omterm{def:b3:forms:closedexact}{tertutup} pada $U$ ($k \geq 1$)
bersifat \omterm{def:b3:forms:closedexact}{eksak}.\index{lema Poincaré}
\end{theorem}

\begin{proof}
Kita bangun sebuah \emph{operator homotopi} linear $h\colon
\Omega^k(U) \to \Omega^{k-1}(U)$ dengan
\begin{equation}\label{eq:b3:forms:homotopy}
\dd(h\omega) + h(\dd\omega) = \omega
\qquad (k \geq 1);
\end{equation}
jika $\dd\omega = 0$, maka $\omega = \dd(h\omega)$ dan selesailah
sudah. Untuk $\omega = \sum_I a_I\,\dd x_I$ tetapkanlah
\[
h\omega = \sum_{\abs I = k}\ \sum_{r=1}^{k}(-1)^{r-1}
\Bigl(\int_0^1 t^{k-1}a_I(tx)\,\dd t\Bigr)\,
x_{i_r}\,\dd x_{i_1}\wedge\dots\wedge
\widehat{\dd x_{i_r}}\wedge\dots\wedge\dd x_{i_k}
\]
(topinya menghapus satu faktor; sedangkan integralnya mulus terhadap $x$
menurut penurunan di bawah tanda integral,
\cref{thm:b3:lebesgue:paramdiff}, sebab semua turunannya
terdominasi pada \omterm{def:b3:topology:compact}{kompak}). Memeriksa
\eqref{eq:b3:forms:homotopy} adalah perhitungan yang dikerjakan sekali
seumur hidup, jadi kita kerjakan selengkapnya. Tetapkan $I$ lalu ambil $\omega =
a\,\dd x_I$ (menurut kelinearannya). Pertama,
\[
\dd(h\omega) = k\Bigl(\int_0^1t^{k-1}a(tx)\dd t\Bigr)\dd x_I
+ \sum_{r=1}^k(-1)^{r-1}\sum_{j=1}^n
\Bigl(\int_0^1t^{k}\,\partial_ja(tx)\,\dd t\Bigr)
x_{i_r}\,\dd x_j\wedge\dd x_{I\setminus i_r} :
\]
kelompok pertamanya mengumpulkan suku yang $\dd$-nya mengenai
faktor $x_{i_r}$ --- sebab \omterm{def:b3:forms:wedge}{hasil kali eksterior} $\dd x_{i_r}\wedge\dd
x_{I\setminus i_r}$ menyusun kembali $\dd x_I$ dengan tanda
$(-1)^{r-1}$ yang meniadakan prafaktornya, dan $k$ nilai
$r$ memberikan faktor $k$ --- sedangkan kelompok keduanya
mengumpulkan suku yang $\dd$-nya mengenai integralnya (sebab aturan
rantainya mengeluarkan $t\,\partial_ja(tx)$). Berikutnya, $\dd\omega =
\sum_j\partial_ja\,\dd x_j\wedge\dd x_I$, dan menerapkan
definisi $h$ pada derajat $k+1$, dengan indeks $j$ menempati
posisi pertamanya:
\[
h(\dd\omega) = \sum_{j=1}^n\Bigl(\int_0^1t^{k}
\partial_ja(tx)\dd t\Bigr)x_j\,\dd x_I
- \sum_{j=1}^n\sum_{r=1}^k(-1)^{r-1}
\Bigl(\int_0^1t^{k}\partial_ja(tx)\dd t\Bigr)
x_{i_r}\,\dd x_j\wedge\dd x_{I\setminus i_r} .
\]
Kedua jumlah gandanya meniadakan diri pada $\dd(h\omega) + h(\dd\omega)$,
yang karena itu sama dengan
\[
\Bigl(\int_0^1\bigl(k\,t^{k-1}a(tx) +
t^k{\textstyle\sum_j}x_j\,\partial_ja(tx)\bigr)\dd t\Bigr)
\dd x_I
= \Bigl(\int_0^1\frac{\dd}{\dd t}\bigl(t^ka(tx)\bigr)\dd
t\Bigr)\dd x_I
= a(x)\,\dd x_I = \omega,
\]
menurut teorema dasar kalkulus. Sifat berbentuk bintangnya
masuk tepat di tempat yang seharusnya: $tx \in U$ bagi $t \in \intcc01$,
sehingga $a(tx)$ bermakna.
\end{proof}

\begin{remark}
Untuk $k = 1$ dan $\omega = \sum_j a_j\dd x_j$, primitifnya
adalah $f(x) = \int_0^1\sum_ja_j(tx)\,x_j\,\dd t$ --- yakni integral
garis $\omega$ sepanjang ruas $[0, x]$: sehingga teoremanya
adalah pernyataan banyak variabel ``medan berjacobi setangkup
adalah gradien'' dari Tahun~2, kini pada setiap derajat. Sedangkan bentuk
sudutnya (\cref{ex:b3:forms:angular}) menunjukkan bahwa hipotesis atas
$U$ bukan hiasan: sebab $\R^2\setminus\{0\}$ tak
berbentuk bintang, dan di sana ketertutupan tak mengakibatkan keeksakan.
Adapun yang bertahan pada \omterm{def:b3:topology:topology}{himpunan terbuka} umum diukur oleh
\emph{kohomologi de Rham} $H^k(U) =
\ker\dd/\operatorname{im}\dd$ --- lihat
\cref{exo:b3:forms:12} untuk perhitungan tak trivial pertamanya.
\end{remark}

\section{Orientasi dan pengintegralan pada submanifold}

Mengintegralkan bentuk-$k$ menuntut wilayah berorientasi
berdimensi $k$. Ingatlah dari \cref{ch:b3:submanifolds}
(\cref{thm:b3:submanifolds:characterizations}) bahwa
submanifold-$k$ $M \subseteq \R^n$ secara lokal merupakan peta sebuah
parametrisasi reguler $\gamma\colon V \to M \cap W$ (dengan $V
\subseteq \R^k$ terbuka, $\gamma$ \omterm{def:b3:topology:continuity}{homeomorfisma} ke petanya
dengan diferensial yang injektif).

\begin{definition}\label{def:b3:forms:orientation}
Sebuah \emph{orientasi}\index{orientasi} $M$ adalah pemilihan,
bagi setiap $p \in M$, salah satu dari dua kelas orientasi
basis \omterm{def:b3:submanifolds:tangent}{ruang singgung} $T_pM$, yang bersifat \emph{koheren
lokal}: bahwa di sekitar setiap titiknya ada parametrisasi
$\gamma$ yang kerangka koordinatnya $(\partial_1\gamma, \dots,
\partial_k\gamma)$ berorientasi positif pada setiap titik
daerah asalnya. Parametrisasi semacam itu disebut \emph{langsung}.
Lalu $M$ disebut \emph{terorientasikan} jika sebuah orientasi ada; dan
pita Möbius menunjukkan bahwa ini bisa gagal. Semua submanifold pada bab
ini berorientasi.
\end{definition}

\begin{definition}[Integral sebuah bentuk]\label{def:b3:forms:integral}
Misalkan $M$ submanifold-$k$ berorientasi dan $\omega$ sebuah
bentuk-$k$ yang terdefinisi pada \omterm{def:b3:topology:topology}{persekitaran} $M$, dengan
$\operatorname{supp}\omega \cap M$ yang \omterm{def:b3:topology:compact}{kompak}.
(a) Jika $\operatorname{supp}\omega \cap M \subseteq
\gamma(V)$ untuk satu parametrisasi langsung, tetapkanlah
\[
\int_M\omega = \int_V\gamma^*\omega
\]
--- dengan ruas kanannya berupa integral Lebesgue atas $V$
(\cref{ch:b3:product}) bagi koefisien $g$ milik
$\gamma^*\omega = g\,\dd u_1\wedge\dots\wedge\dd u_k$, yang
\omterm{def:b3:topology:continuity}{kontinu} dengan tumpuan \omterm{def:b3:topology:compact}{kompak}.
(b) Secara umum, pilihlah berhingga banyak parametrisasi langsung
$\gamma_i(V_i)$ yang menyelimuti \omterm{def:b3:topology:compact}{kompak}
$\operatorname{supp}\omega \cap M$ beserta sebuah
\omterm{lem:b3:forms:partition}{partisi kesatuan} $(\chi_i)$ yang subordinat padanya
(\cref{lem:b3:forms:partition}), lalu tetapkan $\int_M\omega =
\sum_i\int_M\chi_i\,\omega$, dengan setiap sukunya dihitung lewat (a).
Untuk sebuah kurva (dengan $k = 1$) yang diparametrikan oleh
$\gamma\colon\intcc ab\to\R^n$ kita tulis
$\int_\gamma\omega = \int_a^b\gamma^*\omega$, tanpa menuntut
keinjektifan.
\end{definition}

\begin{lemma}[Kekonsistenan]\label{lem:b3:forms:consistency}
Definisi (a) tak bergantung pada parametrisasi
langsungnya, dan definisi (b) tak bergantung baik pada
selimutnya maupun pada \omterm{lem:b3:forms:partition}{partisi kesatuannya}. Lebih jauh, jika
$\Phi$ difeomorfisma antara \omterm{def:b3:topology:topology}{persekitaran} dua submanifold
berorientasi dengan $\Phi(M) = M'$, yang membawa kerangka langsung
$M$ ke kerangka langsung $M'$ pada suatu titik di setiap
komponen $M$, maka $\int_{M'}\omega =
\int_M\Phi^*\omega$.
\end{lemma}

\begin{proof}
(a) Misalkan $\gamma\colon V \to M$, $\delta\colon V' \to M$
parametrisasi langsung yang petanya memuat
$\operatorname{supp}\omega\cap M$. Peralihannya $\tau =
\delta^{-1}\circ\gamma$ merupakan difeomorfisma antara
himpunan bagian terbuka $V, V'$ yang bersangkutan (kemulusan peralihannya:
\cref{thm:b3:submanifolds:characterizations}, lewat pemerian
grafik lokalnya), dan $\gamma = \delta\circ\tau$ di sana, sehingga
$\gamma^*\omega = \tau^*(\delta^*\omega)$
(\cref{thm:b3:forms:pullback}(a)). Tulislah $\delta^*\omega =
g\,\dd u_1\wedge\dots\wedge\dd u_k$; maka
(\cref{thm:b3:forms:pullback}(c)) $\tau^*(\delta^*\omega) =
(g\circ\tau)\,\det(D\tau)\,\dd u_1\wedge\dots\wedge\dd u_k$.
Karena kedua kerangkanya langsung, $D\tau$ membawa basis positif ke
basis positif: $\det D\tau > 0$, sehingga $\det D\tau =
\abs{\det D\tau}$ dan teorema penggantian variabelnya
(\cref{thm:b3:product:changeofvar}) memberikan
$\int(g\circ\tau)\abs{\det D\tau} = \int g$: jadi kedua
integralnya sama. \emph{Inilah seluruh alasan mengapa \omterm{def:b3:forms:orientation}{orientasi}
ada}: bahwa tanpa kendali tandanya, Jacobinya dan nilai
mutlaknya berbeda dan integralnya tak terdefinisi baik.
(b) Jika $(\chi_i)$ dan $(\tilde\chi_j)$ dua partisi yang
sah (termasuk selimutnya), maka menurut (a) dan keaditifan
berhingganya, $\sum_i\int\chi_i\omega =
\sum_{i,j}\int\chi_i\tilde\chi_j\omega =
\sum_j\int\tilde\chi_j\omega$, dengan setiap suku gandanya terhitungkan
pada peta mana pun. Untuk pernyataan terakhirnya: jika $\gamma$ menjelajah
parametrisasi langsung $M$, maka $\Phi\circ\gamma$ menjelajah
parametrisasi langsung $M'$ (sebab perbandingan
\omterm{def:b3:forms:orientation}{orientasinya} konstan lokal, dan tertetapkan pada satu titik per
komponennya), dan $(\Phi\circ\gamma)^*\omega =
\gamma^*(\Phi^*\omega)$.
\end{proof}

\begin{lemma}[Partisi kesatuan, kasus kompak]
\label{lem:b3:forms:partition}
Misalkan $K \subseteq \R^n$ \omterm{def:b3:topology:compact}{kompak} dan $W_1, \dots, W_m$ \omterm{def:b3:topology:topology}{himpunan
terbuka} yang menyelimuti $K$. Maka ada $\chi_1, \dots, \chi_m \in
\mathcal C^\infty(\R^n)$ dengan $0 \leq \chi_i \leq 1$,
$\operatorname{supp}\chi_i \subseteq W_i$ yang \omterm{def:b3:topology:compact}{kompak}, dan
$\sum\chi_i = 1$ pada sebuah \omterm{def:b3:topology:topology}{persekitaran}
$K$.\index{partisi kesatuan}
\end{lemma}

\begin{proof}
Setiap $x \in K$ terletak pada suatu $W_{i(x)}$ dengan bola \omterm{def:b3:forms:closedexact}{tertutup}
$\bar B(x, 2r_x) \subseteq W_{i(x)}$; lalu kekompakannya menyaring
$x_1, \dots, x_N$ dengan bola $B(x_s, r_{x_s})$ yang menyelimuti
$K$. Bagi setiap $s$ ambillah sebuah gundukan $\theta_s \in \mathcal
C^\infty$, $0 \leq \theta_s \leq 1$, $\theta_s = 1$ pada
$\bar B(x_s, r_{x_s})$, $\operatorname{supp}\theta_s
\subseteq B(x_s, 2r_{x_s})$ (muluskanlah indikator
bola berjari-jari $\frac32r_{x_s}$,
\cref{thm:b3:lp:regularization}). Tugaskan setiap $s$ ke satu
indeks $i(s)$ dengan $B(x_s, 2r_{x_s}) \subseteq W_{i(s)}$ lalu
tetapkan $\Theta_i = \sum_{i(s) = i}\theta_s$. Pada \omterm{def:b3:topology:topology}{himpunan terbuka}
$\Omega_0 = \{\sum_j\Theta_j > \tfrac12\} \supseteq K$,
fungsi $\Theta_i/\sum_j\Theta_j$ menyelesaikan tugasnya tetapi hanya
terdefinisi di sana; lalu untuk mengglobalkannya, ambillah $\rho \in \mathcal
C^\infty(\R^n)$ yang memenuhi $\rho = 0$ di tempat $\sum_j\Theta_j
\geq 1$ dan $\rho > 0$ di tempat $\sum_j\Theta_j \leq \tfrac12$
(muluskanlah pemotongan yang sesuai atas $1 - \sum_j\Theta_j$), lalu
tetapkan
\[
\chi_i = \frac{\Theta_i}{\rho + \sum_j\Theta_j} .
\]
Penyebutnya di mana-mana $> 0$ dan sama dengan
$\sum_j\Theta_j$ pada $\{\sum_j\Theta_j \geq 1\}$, yakni sebuah
\omterm{def:b3:topology:topology}{persekitaran} terbuka $K$ (sebab setiap titik $K$ terletak pada suatu bola
$B(x_s, r_{x_s})$ yang $\theta_s = 1$ padanya); dan di sana
$\sum_i\chi_i = 1$. Tumpuan dan batasnya jelas.
\end{proof}

\section{Teorema Stokes}

\begin{definition}\label{def:b3:forms:boundary}
Sebuah \emph{submanifold-$k$ dengan batas}\index{submanifold dengan
batas} $M \subseteq \R^n$ adalah himpunan yang terselimuti
parametrisasi reguler dua jenis: \emph{peta \omterm{def:b3:topology:interior}{interior}}
$\gamma\colon V \to M \cap W$ dengan $V \subseteq \R^k$ terbuka,
dan \emph{peta batas} $\gamma\colon V \cap H^k \to M
\cap W$, dengan $H^k = \{u \in \R^k : u_k \geq 0\}$ dan
$\gamma$ yang diperluas secara mulus dan reguler ke $V$ yang terbuka.
Yang disebut \emph{batas} $\partial M$ adalah himpunan titik yang
tercapai di $u_k = 0$; ia submanifold-$(k-1)$ tanpa
batas, yang diparametrikan oleh pemetaan $u' \mapsto \gamma(u',
0)$. Sebuah \omterm{def:b3:forms:orientation}{orientasi} $M$ menginduksikan satu \omterm{def:b3:forms:orientation}{orientasi} pada $\partial M$ lewat
aturan \emph{normal-keluar-dulu}: bahwa di $p \in \partial M$,
sebuah basis $(w_1, \dots, w_{k-1})$ milik $T_p\partial M$ bersifat
positif jika dan hanya jika $(\nu, w_1, \dots, w_{k-1})$ merupakan basis
positif $T_pM$, dengan $\nu \in T_pM \setminus T_p\partial M$
yang menunjuk ke luar $M$ (dan pada peta batasnya: $\nu =
-\partial_k\gamma$, sampai penambahan komponen singgungnya ---
sebab kelas \omterm{def:b3:forms:orientation}{orientasinya} tak melihatnya).
\end{definition}

\begin{omfigure}
\begin{tikzpicture}[scale=1.05]
  \fill[omDef!10] (-2,0) arc (180:0:2) -- cycle;
  \draw[thick, omDef] (-2,0) arc (180:0:2);
  \draw[thick, omThm] (-2,0) -- (2,0);
  \node[omDef] at (0.1,1.35) {$M$};
  \node[omThm, below] at (1.55,-0.05) {$\partial M$};
  \draw[->, thick, omExo] (1.415,1.415) -- (2.05,2.05);
  \node[omExo, right] at (2.02,2.12) {$\nu$};
  \draw[->, thick] (1.29,1.53) arc (50:75:2);
  \draw[->, thick, omExo] (-0.6,0) -- (-0.6,-0.65);
  \node[omExo, below] at (-0.6,-0.65) {$\nu$};
  \draw[->, thick] (-0.35,0) -- (0.55,0);
  \node[below] at (0.1,-0.14) {terinduksi};
\end{tikzpicture}

{\small Aturan normal-keluar-dulu: bahwa pada setiap titik
batasnya, taruhlah vektor keluar $\nu$ lebih dahulu; lalu basis yang
melengkapinya menjadi kerangka positif $M$ mengorientasikan $\partial M$.
Untuk daerah bidang dengan \omterm{def:b3:forms:orientation}{orientasi} bakunya, inilah aturan
berlawanan arah jarum jam milik Green--Riemann.}
\end{omfigure}

\begin{lemma}[Stokes pada setengah ruang]
\label{lem:b3:forms:halfspace}
Misalkan $\eta$ bentuk-$(k-1)$ mulus pada $\R^k$ dengan tumpuan
\omterm{def:b3:topology:compact}{kompak}. Maka
\[
\int_{H^k}\dd\eta = \int_{\partial H^k}\eta,
\]
dengan $H^k$ mengemban \omterm{def:b3:forms:orientation}{orientasi} baku $\R^k$ dan
$\partial H^k = \{u_k = 0\} \cong \R^{k-1}$ mengemban \omterm{def:b3:forms:orientation}{orientasi} terinduksinya,
yang merupakan $(-1)^k$ kali \omterm{def:b3:forms:orientation}{orientasi} baku
$\R^{k-1}$.
\end{lemma}

\begin{proof}
Pertama pembukuan \omterm{def:b3:forms:orientation}{orientasinya}: normal keluar sepanjang
$\partial H^k$ adalah $-e_k$, dan
\[
\det(-e_k, e_1, \dots, e_{k-1}) =
-\det(e_k, e_1, \dots, e_{k-1}) = -(-1)^{k-1}\det(e_1,
\dots, e_k) = (-1)^k :
\]
kerangka $(e_1, \dots, e_{k-1})$ milik $\partial H^k$ bersifat
positif bagi \omterm{def:b3:forms:orientation}{orientasi} terinduksinya tepat saat $k$
genap, dan dari sanalah perbandingan yang dinyatakan itu. Menurut kelinearannya ambillah $\eta
= f\,\dd u_1\wedge\dots\wedge\widehat{\dd
u_i}\wedge\dots\wedge\dd u_k$, $f \in \mathcal
C^\infty_c(\R^k)$; maka $\dd\eta =
(-1)^{i-1}\,\partial_if\,\dd u_1\wedge\dots\wedge\dd u_k$
(sebab memindahkan $\dd u_i$ ke posisinya memakan $i - 1$ penukaran). Ada dua
kasus, keduanya lewat Tonelli--Fubini
(\cref{thm:b3:product:fubini}) dan teorema dasar kalkulus
satu variabel.

\emph{Kasus $i < k$.} Dengan mengintegralkan lebih dahulu terhadap $u_i$ atas $\R$:
$\int_\R\partial_if\,\dd u_i = 0$ (menurut tumpuan \omterm{def:b3:topology:compact}{kompaknya}), sehingga
$\int_{H^k}\dd\eta = 0$. Sedangkan pembatasan $\eta$ ke
$\{u_k = 0\}$ memuat faktor $\dd u_k$, yang terbatas
menjadi $0$ (sebab $u_k$ konstan di sana): sehingga $\int_{\partial H^k}\eta =
0$ juga.

\emph{Kasus $i = k$.} Dengan mengintegralkan lebih dahulu terhadap $u_k$ atas
$\intco0\infty$:
\begin{align*}
\int_{H^k}\dd\eta &= (-1)^{k-1}\int_{\R^{k-1}}
\Bigl(\int_0^\infty\partial_kf\,\dd u_k\Bigr)\dd u' \\
&= (-1)^{k-1}\int_{\R^{k-1}}\bigl(0 - f(u', 0)\bigr)\dd u'
= (-1)^k\int_{\R^{k-1}}f(u',0)\,\dd u' .
\end{align*}
Pada sisi batasnya, $\eta$ terbatas menjadi $f(u',
0)\,\dd u_1\wedge\dots\wedge\dd u_{k-1}$, dan karena \omterm{def:b3:forms:orientation}{orientasi}
terinduksinya $(-1)^k$ kali \omterm{def:b3:forms:orientation}{orientasi} bakunya,
$\int_{\partial H^k}\eta =
(-1)^k\int_{\R^{k-1}}f(u',0)\,\dd u'$. Jadi kedua ruasnya sama.
\end{proof}

\begin{theorem}[Stokes]\label{thm:b3:forms:stokes}
Misalkan $M \subseteq \R^n$ submanifold-$k$ \omterm{def:b3:topology:compact}{kompak} berorientasi
dengan \omterm{def:b3:forms:boundary}{batas}, dengan $\partial M$ mengemban \omterm{def:b3:forms:orientation}{orientasi}
terinduksinya, dan misalkan $\omega$ bentuk-$(k-1)$ mulus pada sebuah
\omterm{def:b3:topology:topology}{persekitaran} $M$. Maka\index{teorema Stokes}
\[
\int_M\dd\omega = \int_{\partial M}\omega .
\]
Khususnya, jika $\partial M = \varnothing$:
$\int_M\dd\omega = 0$.
\end{theorem}

\begin{proof}
Selimutilah \omterm{def:b3:topology:compact}{kompak} $M$ dengan berhingga banyak peta langsung
(\omterm{def:b3:topology:interior}{interior} atau \omterm{def:b3:forms:boundary}{batas}), lalu ambillah \omterm{lem:b3:forms:partition}{partisi kesatuan}
$(\chi_i)$ yang subordinat pada \omterm{def:b3:topology:topology}{himpunan terbuka} $W_i$ yang bersesuaian
di $\R^n$ (\cref{lem:b3:forms:partition}), dengan $\sum\chi_i
= 1$ pada sebuah \omterm{def:b3:topology:topology}{persekitaran} $M$. Pada \omterm{def:b3:topology:topology}{persekitaran} itu
$\dd(\sum\chi_i) = 0$, sehingga
\[
\int_M\dd\omega = \sum_i\int_M\dd(\chi_i\omega),
\qquad
\int_{\partial M}\omega =
\sum_i\int_{\partial M}\chi_i\omega :
\]
kedua ruasnya aditif, dan cukup dibuktikan
teoremanya bagi bentuk yang bertumpu pada satu peta saja.

\emph{Peta \omterm{def:b3:topology:interior}{interior}.} Jika $\operatorname{supp}\omega \cap M
\subseteq \gamma(V)$ dengan $V$ terbuka di $\R^k$: perluaslah $\eta =
\gamma^*\omega$ oleh nol ke $\R^k$ (mulus, bertumpuan \omterm{def:b3:topology:compact}{kompak}
di $V$) lalu terapkan \cref{lem:b3:forms:halfspace} dengan
tumpuannya menjauh dari $\partial H^k$ (geserlah $V$ ke
setengah ruang atas yang terbuka --- atau ulangilah saja perhitungan kasus $i < k$
atas seluruh $\R^k$): $\int_M\dd\omega =
\int_{\R^k}\dd\eta = 0$, dan $\int_{\partial M}\omega = 0$
sebab $\omega$ lenyap di dekat $\partial M$.

\emph{\omterm{def:b3:forms:boundary}{Peta batas}.} Jika $\operatorname{supp}\omega\cap M
\subseteq \gamma(V \cap H^k)$: dengan $\eta = \gamma^*\omega$
yang diperluas oleh nol,
$\gamma^*(\dd\omega) = \dd\eta$
(\cref{thm:b3:forms:pullback}(b)), sehingga menurut
\cref{def:b3:forms:integral}:
\[
\int_M\dd\omega = \int_{H^k}\dd\eta
\overset{\text{\cref{lem:b3:forms:halfspace}}}{=}
\int_{\partial H^k}\eta .
\]
Tinggal mengenali ruas kanannya dengan $\int_{\partial
M}\omega$. Batasnya diparametrikan oleh $\beta(u') =
\gamma(u', 0)$, dan $\beta^*\omega$ merupakan pembatasan
$\eta$ ke $\{u_k = 0\}$ (yakni \omterm{def:b3:forms:pullback}{tarikan balik} di bawah inklusi $u'
\mapsto (u', 0)$ yang dikomposisikan dengan $\gamma$). Perbandingan
\omterm{def:b3:forms:orientation}{orientasinya} adalah $(-1)^k$ yang sama pada kedua ruasnya: sebab kerangka
$(\partial_1\beta, \dots, \partial_{k-1}\beta)$ duduk pada
\omterm{def:b3:forms:orientation}{orientasi} terinduksi $\partial M$ dengan tanda
$\det(-e_k, e_1, \dots, e_{k-1}) = (-1)^k$ yang dihitung pada
petanya (sebab vektor keluarnya tertarik balik menjadi $-e_k$), yang persis
merupakan tanda yang mengaitkan \omterm{def:b3:forms:orientation}{orientasi} terinduksi $\partial H^k$
dengan $\R^{k-1}$ yang baku
(\cref{lem:b3:forms:halfspace}). Jadi kedua konvensi tandanya
meniadakan diri:
$\int_{\partial H^k}\eta = \int_{\partial M}\omega$.
\end{proof}

\begin{example}[Teorema klasiknya]
\label{ex:b3:forms:classical}
Misalkan $D \subseteq \R^2$ daerah \omterm{def:b3:topology:compact}{kompak} dengan \omterm{def:b3:forms:boundary}{batas}
mulus, yang terorientasi secara baku. Untuk $\omega = P\,\dd x +
Q\,\dd y$: $\dd\omega = \bigl(\partial_xQ -
\partial_yP\bigr)\dd x\wedge\dd y$, dan Stokes berbunyi
\[
\int_D\Bigl(\frac{\partial Q}{\partial x} -
\frac{\partial P}{\partial y}\Bigr)\dd x\,\dd y
= \oint_{\partial D}P\,\dd x + Q\,\dd y :
\]
yakni \emph{Green--Riemann}, yang dibuktikan bagi daerah elementer pada
jilid Tahun~2 dan kini dalam keumuman alaminya. Di $\R^3$,
Stokes yang diterapkan pada bentuk-$2$ fluks sebuah medan vektor pada
daerah \omterm{def:b3:topology:compact}{kompak} memberikan \emph{teorema divergensi}
$\int_\Omega\operatorname{div}F =
\int_{\partial\Omega}\langle F, \nu\rangle\,\dd S$, sedangkan
yang diterapkan pada bentuk-$1$ pada permukaan-dengan-batas memberikan
teorema curl \emph{Kelvin--Stokes} yang klasik;
\cref{exo:b3:forms:7} menguraikan kedua kamusnya.
\end{example}

\section{Bilangan lilitan}\label{sec:b3:forms:winding}

\begin{definition}\label{def:b3:forms:winding}
Misalkan $\gamma\colon\intcc01\to\R^2\setminus\{a\}$ kurva
\omterm{def:b3:forms:closedexact}{tertutup} yang mulus. Yang disebut \emph{bilangan lilitan}\index{bilangan
lilitan} kurva itu di sekitar $a$ adalah
\[
\operatorname{Ind}_\gamma(a) =
\frac1{2\pi}\int_\gamma\omega_\theta^a,
\qquad
\omega_\theta^a = \frac{(x - a_1)\,\dd y - (y - a_2)\,\dd
x}{(x - a_1)^2 + (y - a_2)^2} .
\]
\end{definition}

\begin{proposition}\label{prop:b3:forms:windinginteger}
Berlaku $\operatorname{Ind}_\gamma(a) \in \Z$; dan sebagai fungsi $a$
ia konstan pada setiap \omterm{def:b3:topology:components}{komponen terhubung}
$\R^2\setminus\gamma(\intcc01)$ serta nol pada komponen
tak terbatasnya. Untuk $\gamma(t) = a + r(\cos2\pi t, \sin2\pi t)$:
$\operatorname{Ind}_\gamma(a) = 1$.
\end{proposition}

\begin{proof}
Ambillah $a = 0$ lalu tulis $\gamma = (x, y)$, $\rho =
\norm\gamma > 0$, $\vartheta(t) =
\int_0^t\gamma^*\omega_\theta$, sehingga $\vartheta' =
\frac{xy' - yx'}{x^2 + y^2}$. Dalam notasi kompleksnya ambillah $u(t)
= \gamma(t)\,\rho(t)^{-1}\eu^{-\iu\vartheta(t)}$; maka
$\abs u = 1$ dan perhitungan langsung memberikan
\[
\frac{u'}{u} = \frac{\gamma'}{\gamma} - \frac{\rho'}{\rho}
- \iu\vartheta'
= \frac{\bar\gamma\gamma'}{\abs\gamma^2} -
\frac{\rho'}{\rho} - \iu\,\frac{xy' -
yx'}{\rho^2}
= \frac{(xx' + yy')}{\rho^2} - \frac{\rho'}{\rho} = 0,
\]
sebab $\bar\gamma\gamma' = (xx' + yy') + \iu(xy' - yx')$ dan
$\rho\rho' = xx' + yy'$. Jadi $u$ konstan: $\gamma(t) =
\rho(t)\,c\,\eu^{\iu\vartheta(t)}$ dengan $\abs c = 1$, sedangkan
$\gamma(1) = \gamma(0)$ beserta $\rho(1) = \rho(0)$ memaksa
$\eu^{\iu\vartheta(1)} = \eu^{\iu\vartheta(0)} = 1$:
$\vartheta(1) \in 2\pi\Z$, yakni
$\operatorname{Ind}_\gamma(0) \in \Z$. Lalu sebagai fungsi $a$
pada komplemen terbuka kurva \omterm{def:b3:topology:compact}{kompaknya}, integral
pendefinisinya \omterm{def:b3:topology:continuity}{kontinu} (\cref{thm:b3:lebesgue:paramcont},
lewat dominasi pada \omterm{def:b3:topology:topology}{persekitaran} setiap $a$); sedangkan fungsi \omterm{def:b3:topology:continuity}{kontinu}
bernilai bulat bersifat konstan lokal, sehingga konstan
pada komponennya. Untuk $\norm a$ yang besar integrannya
$O(1/\norm a)$ seragam terhadap $t$, sehingga indeksnya menuju $0$
dan lenyap pada komponen tak terbatasnya. Untuk lingkarannya:
$\gamma^*\omega_\theta^a = 2\pi\,\dd t$ secara langsung.
\end{proof}

\begin{remark}
Di bawah $\C \cong \R^2$, $\frac{\dd z}{z} = \frac{x\,\dd x +
y\,\dd y}{x^2 + y^2} + \iu\,\omega_\theta$, sehingga
$\operatorname{Ind}_\gamma(a) =
\frac1{2\iu\pi}\oint_\gamma\frac{\dd z}{z - a}$: inilah
indeks \cref{ch:b3:residues}, dan
\cref{prop:b3:forms:windinginteger} membuktikan ulang
kebulatan dan kekonstanan lokalnya lewat cara variabel real ---
yakni separuh \omterm{def:b3:topology:topology}{topologis} teorema \omterm{def:b3:residues:singularities}{residunya}, kini berdiri di atas
Stokes.
\end{remark}

\begin{proposition}\label{prop:b3:forms:exactloop}
Jika $\omega = \dd f$ \omterm{def:b3:forms:closedexact}{eksak} pada $U$ yang terbuka dan
$\gamma\colon\intcc01\to U$ kurva \omterm{def:b3:forms:closedexact}{tertutup}, maka
$\int_\gamma\omega = 0$.
\end{proposition}

\begin{proof}
$\int_\gamma\dd f = \int_0^1(f\circ\gamma)'(t)\,\dd t =
f(\gamma(1)) - f(\gamma(0)) = 0$ --- sebab aturan rantainya
mengenali $\gamma^*(\dd f)$ dengan $(f\circ\gamma)'\,\dd t$.
\end{proof}

\begin{method}\label{met:b3:forms:method}
Menghitung dengan bentuk: (1) mekanisasikanlah --- \omterm{def:b3:forms:wedge}{hasil kali eksteriornya}
diurutkan ulang dengan tanda, $\dd$ menurunkan koefisiennya menjadi $\dd
x_j$ yang baru, dan \omterm{def:b3:forms:pullback}{tarikan baliknya} mensubstitusi; percayailah aljabarnya, sebab ia menyandikan
setiap Jacobi. (2) Untuk mengintegralkan sebuah bentuk atas submanifold:
parametrikanlah secara langsung, tariklah balik, lalu integralkan koefisiennya;
\omterm{def:b3:forms:orientation}{orientasinya} satu-satunya jebakan --- periksalah satu kerangka. (3) Untuk
membuktikan sebuah identitas integral, carilah rupa Stokes: apakah
integrannya \omterm{def:b3:forms:closedexact}{eksak}? apakah daerahnya sebuah \omterm{def:b3:forms:boundary}{batas}? (4) Untuk membandingkan
integral atas dua submanifold yang ``sejajar'', terapkanlah Stokes
pada daerah di antara keduanya (yakni argumen deformasinya,
\cref{exo:b3:forms:10}). (5) Integral tak nol sebuah \omterm{def:b3:forms:closedexact}{bentuk
tertutup} mengesahkan adanya halangan \omterm{def:b3:topology:topology}{topologis} --- tak ada primitif,
tak ada retraksi, tak ada perluasan bebas nol: dan beginilah
soal akhir pekannya membunuh retraksi bolanya.
\end{method}

\section{Latihan}

\begin{exercise}[$\star$]\label{exo:b3:forms:1}
Pada $\R^3$, misalkan $\omega = x\,\dd y - z\,\dd x$ dan $\eta =
\dd x + y\,\dd z$. Hitunglah $\omega\wedge\eta$, $\dd\omega$,
$\dd\eta$, dan $\dd(\omega\wedge\eta)$, lalu periksalah
aturan Leibniz berderajatnya pada contoh ini.
\end{exercise}

\begin{exercise}[$\star$]\label{exo:b3:forms:2}
Kenalilah, pada $\R^3$, ketiga penjelmaan $\dd$: bagi
$f \in \Omega^0$, $\dd f \leftrightarrow \nabla f$; bagi
bentuk kerja $\omega_F = F_1\dd x + F_2\dd y + F_3\dd z$,
$\dd\omega_F \leftrightarrow \operatorname{curl}F$; dan bagi
bentuk fluks $\sigma_F = F_1\,\dd y\wedge\dd z + F_2\,\dd
z\wedge\dd x + F_3\,\dd x\wedge\dd y$, $\dd\sigma_F
\leftrightarrow \operatorname{div}F$. Simpulkanlah dari $\dd^2 =
0$ identitas $\operatorname{curl}\nabla f = 0$ dan
$\operatorname{div}\operatorname{curl}F = 0$.
\end{exercise}

\begin{exercise}[$\star\star$]\label{exo:b3:forms:3}
Putuskanlah apakah setiap bentuk-$1$ berikut \omterm{def:b3:forms:closedexact}{tertutup}, \omterm{def:b3:forms:closedexact}{eksak} pada daerah asalnya,
lalu hitunglah primitifnya bila ada:
(a) $(2xy + z^2)\,\dd x + x^2\,\dd y + 2xz\,\dd z$ pada
$\R^3$;
(b) $\dfrac{x\,\dd x + y\,\dd y}{x^2 + y^2}$ pada
$\R^2\setminus\{0\}$;
(c) $\dfrac{-y\,\dd x + x\,\dd y}{x^2 + y^2}$ pada
setengah bidang $\{x > 0\}$.
\end{exercise}

\begin{exercise}[$\star\star$]\label{exo:b3:forms:4}
(Bentuk sudutnya) Periksalah bahwa $\omega_\theta$
(\cref{ex:b3:forms:angular}) \omterm{def:b3:forms:closedexact}{tertutup}; hitunglah
$\int_\gamma\omega_\theta$ bagi $\gamma$ lingkaran berjari-jari
$r$ di sekitar $0$; lalu simpulkan bahwa $\omega_\theta$ tak \omterm{def:b3:forms:closedexact}{eksak}
pada $\R^2\setminus\{0\}$, dan bahwa $\R^2\setminus\{0\}$ tak
berbentuk bintang terhadap satu titik pun yang dimilikinya (lewat dua jalur:
lewat \cref{thm:b3:forms:poincare}, dan langsung dari
geometrinya).
\end{exercise}

\begin{exercise}[$\star\star$]\label{exo:b3:forms:5}
(Bentuk luas sebuah hiperpermukaan) Misalkan $M \subseteq \R^n$
hiperpermukaan \omterm{def:b3:topology:compact}{kompak} berorientasi yang \omterm{def:b3:forms:orientation}{orientasinya} diberikan oleh
medan normal satuan $\nu$ (dengan $(w_1, \dots, w_{n-1})$ positif
jika dan hanya jika $(\nu, w_1, \dots, w_{n-1})$ positif di $\R^n$). Tunjukkanlah
bahwa bentuk-$(n-1)$ $\sigma_\nu =
\sum_i(-1)^{i-1}\nu_i\,\dd x_1\wedge\dots\wedge\widehat{\dd
x_i}\wedge\dots\wedge\dd x_n$ terbatas pada $M$ menjadi
\emph{bentuk luasnya}: bahwa bagi parametrisasi langsung $\gamma$,
$\gamma^*\sigma_\nu = \sqrt{\det G}\,\dd u_1\wedge\dots
\wedge\dd u_{n-1}$, dengan $G = ({}^t D\gamma)(D\gamma)$ sebagai
matriks Gram. \emph{(Perhatikanlah bahwa
$(\gamma^*\sigma_\nu)(e_1, \dots, e_{n-1}) = \det(\nu,
\partial_1\gamma, \dots, \partial_{n-1}\gamma)$ lalu kuadratkanlah
determinan itu.)} Hitunglah $\int_{S^2}\sigma_\nu = 4\pi$.
\end{exercise}

\begin{exercise}[$\star\star$]\label{exo:b3:forms:6}
Hitunglah $\int_{S^2}\omega$ bagi $\omega = x\,\dd y\wedge\dd z
+ y\,\dd z\wedge\dd x + z\,\dd x\wedge\dd y$: (a) langsung
dalam koordinat bolanya; (b) lewat Stokes pada bola satuannya.
Simpulkanlah $\operatorname{vol}(B^3) = \frac{4\pi}3$ dari
luas $S^2$, lalu umumkanlah:
$n\operatorname{vol}(B^n) = \operatorname{area}(S^{n-1})$,
yang selaras dengan \cref{thm:b3:product:ballvolume}.
\end{exercise}

\begin{exercise}[$\star\star$]\label{exo:b3:forms:7}
(Kamusnya) Turunkanlah dengan cermat dari
\cref{thm:b3:forms:stokes}: (a) teorema divergensi di
$\R^3$ (gabungkanlah \cref{exo:b3:forms:2} dan
\cref{exo:b3:forms:5}); (b) teorema Kelvin--Stokes
$\int_S\langle\operatorname{curl}F, \nu\rangle\,\dd S =
\oint_{\partial S}\langle F, \tau\rangle\,\dd\ell$ bagi
permukaan \omterm{def:b3:topology:compact}{kompak} berorientasi dengan \omterm{def:b3:forms:boundary}{batas} di $\R^3$. Periksalah bahwa
konvensi \omterm{def:b3:forms:orientation}{orientasinya} cocok pada setengah bola atas yang
dibatasi khatulistiwanya.
\end{exercise}

\begin{exercise}[$\star\star$]\label{exo:b3:forms:8}
(Identitas Green) Untuk $u, v$ yang mulus pada \omterm{def:b3:topology:topology}{persekitaran}
daerah \omterm{def:b3:topology:compact}{kompak} $\Omega \subseteq \R^n$ yang berbatas
mulus, buktikanlah
\[
\int_\Omega\bigl(u\,\Delta v + \langle\nabla u, \nabla
v\rangle\bigr) = \int_{\partial\Omega}u\,\partial_\nu
v\,\dd S,
\qquad
\int_\Omega\bigl(u\,\Delta v - v\,\Delta u\bigr) =
\int_{\partial\Omega}\bigl(u\,\partial_\nu v -
v\,\partial_\nu u\bigr)\dd S .
\]
Simpulkanlah: bahwa \omterm{prop:b3:conformal:harmonicholo}{fungsi harmonik} pada $\Omega$ yang lenyap pada
$\partial\Omega$ lenyap secara identik, dan dua \omterm{prop:b3:conformal:harmonicholo}{fungsi
harmonik} dengan nilai \omterm{def:b3:forms:boundary}{batas} yang sama berimpit ---
yakni ketunggalan pada masalah Dirichlet
\cref{ch:b3:conformal}.
\end{exercise}

\begin{exercise}[$\star\star$]\label{exo:b3:forms:9}
(Penggantian variabel, bentuk berorientasi) Misalkan $\varphi\colon U
\to V$ difeomorfisma antara \omterm{def:b3:topology:topology}{himpunan terbuka} $\R^n$ dengan $\det
D\varphi > 0$, dan $f$ \omterm{def:b3:topology:continuity}{kontinu} bertumpuan \omterm{def:b3:topology:compact}{kompak} di
$V$. Tunjukkanlah bahwa identitas \omterm{def:b3:forms:pullback}{tarikan baliknya}
$\int_U\varphi^*(f\,\dd x_1\wedge\dots\wedge\dd x_n) =
\int_Vf\,\dd x_1\wedge\dots\wedge\dd x_n$ setara dengan
teorema penggantian variabelnya
(\cref{thm:b3:product:changeofvar}) bagi $\varphi$ semacam itu, lalu
jelaskanlah persis ke mana nilai mutlak pada Jacobinya
pergi.
\end{exercise}

\begin{exercise}[$\star\star\star$]\label{exo:b3:forms:10}
(Deformasi) Misalkan $\omega$ bentuk-$2$ \emph{\omterm{def:b3:forms:closedexact}{tertutup}} pada
$\R^3\setminus\{0\}$ dan $S_r$ bola berjari-jari $r$
yang berpusat di $0$. Tunjukkanlah bahwa $\int_{S_r}\omega$ tak
bergantung pada $r > 0$ \emph{(terapkanlah Stokes pada cangkang antara
dua jari-jari; perhatikanlah kedua \omterm{def:b3:forms:orientation}{orientasi} batasnya)}. Terapkanlah pada
\emph{bentuk sudut ruang}
\[
\omega = \frac{x\,\dd y\wedge\dd z + y\,\dd z\wedge\dd x +
z\,\dd x\wedge\dd y}{(x^2 + y^2 + z^2)^{3/2}} :
\]
periksalah bahwa ia \omterm{def:b3:forms:closedexact}{tertutup}, hitunglah $\int_{S_r}\omega = 4\pi$, lalu
simpulkanlah bahwa ia \omterm{def:b3:forms:closedexact}{tertutup} tetapi tak \omterm{def:b3:forms:closedexact}{eksak} pada
$\R^3\setminus\{0\}$ --- yakni saudara berdimensi dua
$\omega_\theta$, dan isi geometris hukum Gauss
dalam elektrostatika.
\end{exercise}

\begin{exercise}[$\star\star$]\label{exo:b3:forms:11}
Misalkan $\gamma$ kurva \omterm{def:b3:forms:closedexact}{tertutup} mulus di
$\R^2\setminus\{0\}$ dengan $n = \operatorname{Ind}_\gamma(0)$.
Tunjukkanlah $\int_\gamma\omega = n\int_{S^1}\omega$ bagi
\emph{setiap} bentuk-$1$ \omterm{def:b3:forms:closedexact}{tertutup} $\omega$ pada
$\R^2\setminus\{0\}$ \emph{(tulislah $\omega = c\,\omega_\theta
+ \dd f$ menurut \cref{exo:b3:forms:12})}. Tafsirannya: bahwa pada
bidang tertusuknya, \omterm{def:b3:forms:winding}{bilangan lilitan} adalah satu-satunya halangan
bagi lenyapnya periode.
\end{exercise}

\begin{exercise}[$\star\star\star$]\label{exo:b3:forms:12}
(Perhitungan de Rham pertama) Tunjukkanlah bahwa setiap bentuk-$1$ \omterm{def:b3:forms:closedexact}{tertutup}
$\omega$ pada $U = \R^2\setminus\{0\}$ secara tunggal berupa
\[
\omega = c\,\omega_\theta + \dd f,
\qquad c = \frac1{2\pi}\int_{S^1}\omega,\quad f \in
\mathcal C^\infty(U) :
\]
definisikanlah $f(p)$ dengan mengintegralkan $\omega - c\,\omega_\theta$
sepanjang \omterm{def:b3:topology:pathconnected}{lintasan} dari $(1,0)$ ke $p$ (potongan radial lalu busur
lingkaran), tunjukkanlah bahwa hasilnya tak bergantung pada pilihannya
justru karena periode-$S^1$-nya lenyap, lalu periksalah $\dd
f = \omega - c\,\omega_\theta$. Simpulkanlah:
$H^1(\R^2\setminus\{0\}) \cong \R$, yang dibangkitkan bentuk
sudutnya.
\end{exercise}

\section{Soal: teorema titik tetap Brouwer}

\begin{problem}[{Soal akhir pekan --- tak ada retraksi, tak ada
pelolosan}]\label{pb:b3:forms:1}
Teorema Brouwer menyatakan bahwa \emph{setiap \omterm{def:b3:topology:continuity}{pemetaan kontinu} dari
bola satuan \omterm{def:b3:forms:closedexact}{tertutup} $\bar B = \bar B^n \subseteq \R^n$ ke
dirinya sendiri mempunyai titik tetap} --- yakni salah satu teorema besar
matematika, dengan akibat mulai dari teori permainan (kesetimbangan
Nash) sampai analisis matriks. Buktinya lewat \omterm{def:b3:forms:diffform}{bentuk diferensial}
adalah yang terbersih yang dikenal: Stokes menunjukkan bahwa sferanya bukan
retrak bolanya, dan selebihnya menyusul. Sepanjang soal ini,
$S = S^{n-1} = \partial\bar B$, $n \geq 2$, dan
\[
\sigma = \sum_{i=1}^n(-1)^{i-1}x_i\,
\dd x_1\wedge\dots\wedge\widehat{\dd x_i}\wedge\dots\wedge
\dd x_n
\]
(topinya menghapus faktornya).

\medskip
\noindent\textbf{Bagian I --- Alat ukurnya.}
\begin{enumerate}
  \item Hitunglah $\dd\sigma$, lalu simpulkanlah dari Stokes
        (\cref{thm:b3:forms:stokes}) bahwa $\int_S\sigma =
        n\operatorname{vol}(\bar B) > 0$, dengan $S$ mengemban
        \omterm{def:b3:forms:orientation}{orientasi} \omterm{def:b3:forms:boundary}{batas} bolanya.
  \item Untuk $n = 2$ dan $n = 3$, kenalilah pembatasan
        $\sigma$ ke $S$ sebagai bentuk panjang busur dan bentuk luasnya
        (\cref{exo:b3:forms:5} dengan $\nu(x) = x$) lalu
        hitunglah ulang $\int_S\sigma$ secara langsung.
  \item Misalkan $W \subseteq \R^N$ terbuka dan $\varphi\colon W
        \to \R^n$ mulus dengan $\norm{\varphi(x)} = 1$ bagi
        setiap $x \in W$. Tunjukkanlah bahwa
        $\varphi^*(\dd\sigma) = 0$. \emph{(Turunkanlah
        $\norm\varphi^2 = 1$: peta $D\varphi(x)$
        terletak pada hiperbidang $\varphi(x)^\perp$, yang
        berdimensi $n - 1$; lalu terapkanlah
        \cref{prop:b3:forms:linearpullback}(2).)}
  \item Di manakah cacat ``bukti'' berikut bahwa
        $\int_S\sigma = 0$: ``$S$ \omterm{def:b3:topology:compact}{kompak} tanpa
        \omterm{def:b3:forms:boundary}{batas}, dan $\sigma$ yang dibatasi ke $S$ merupakan bentuk
        berderajat tertinggi padanya, sehingga \omterm{def:b3:forms:closedexact}{tertutup}, sehingga $\int_S\sigma =
        \int_S\dd(\text{sesuatu}) = 0$ menurut Stokes''?
        (Tunjuklah kata yang keliru.)
  \item Jelaskanlah dalam satu paragraf strategi Bagian II:
        apa yang akan diintegralkan, atas apa, dan dari mana
        kontradiksinya akan datang.
\end{enumerate}

\medskip
\noindent\textbf{Bagian II --- Tak ada retraksi mulus.} Andaikan,
demi kontradiksi, bahwa $r$ merupakan \emph{retraksi mulus} dari
bolanya ke sferanya: $r$ mulus pada \omterm{def:b3:topology:topology}{persekitaran}
$\bar B$, $r(\bar B) \subseteq S$, dan $r(x) = x$ bagi setiap
$x \in S$.
\begin{enumerate}[resume]
  \item Benarkanlah $\int_Sr^*\sigma = \int_S\sigma$
        \emph{(pada $S$, $r$ terbatas menjadi identitasnya: sebab jika
        $\gamma$ parametrisasi langsung sekeping
        $S$, maka $r\circ\gamma = \gamma$)}.
  \item Dengan memakai Stokes pada $\bar B$, tunjukkanlah $\int_Sr^*\sigma =
        \int_{\bar B}\dd(r^*\sigma)$.
  \item Tunjukkanlah $\dd(r^*\sigma) = r^*(\dd\sigma) = 0$ (yakni pertanyaan
        3 yang diterapkan pada $\varphi = r$), lalu simpulkanlah:
        \emph{tak ada retraksi mulus $\bar B \to S$.}
  \item Selesaikanlah kasus terkecuali $n = 1$ dengan tangan: tunjukkanlah
        secara langsung bahwa tak ada \omterm{def:b3:topology:continuity}{pemetaan kontinu} $\intcc{-1}1 \to
        \{-1, 1\}$ yang menetapkan kedua ujungnya, lalu sebutkanlah
        teorema yang dipakai.
\end{enumerate}

\medskip
\noindent\textbf{Bagian III --- Brouwer mulus.} Misalkan $g$
mulus pada \omterm{def:b3:topology:topology}{persekitaran} $\bar B$ dengan $g(\bar B)
\subseteq \bar B$ dan \emph{tanpa} titik tetap di $\bar B$.
\begin{enumerate}[resume]
  \item Tunjukkanlah $\delta = \min_{x\in\bar B}\norm{g(x) - x} >
        0$.
  \item Untuk $x \in \bar B$ ambillah $u(x) = \frac{x -
        g(x)}{\norm{x - g(x)}}$ dan misalkan $r(x) = x +
        t(x)\,u(x)$ perpotongan sinar
        $\{x + tu(x) : t \geq 0\}$ dengan $S$. Selesaikanlah
        persamaan kuadrat $\norm{x + tu}^2 = 1$ lalu perolehlah
        \[
        t(x) = -\langle x, u(x)\rangle +
        \sqrt{1 - \norm x^2 + \langle x, u(x)\rangle^2}
        \;\geq\; 0 .
        \]
  \item Tunjukkanlah bahwa isi akarnya positif \emph{sejati}
        pada $\bar B$: sebab jika $1 - \norm x^2 + \langle x,
        u(x)\rangle^2 = 0$ maka $\norm x = 1$ dan $\langle
        x, u(x)\rangle = 0$, yakni $\langle x, x -
        g(x)\rangle = 0$, yakni $\langle x, g(x)\rangle =
        1$; lalu menurut Cauchy--Schwarz dengan $\norm x = 1$,
        $\norm{g(x)} \leq 1$, ini memaksa $g(x) = x$ ---
        yang terkecualikan. Simpulkanlah bahwa $r$ mulus pada
        \omterm{def:b3:topology:topology}{persekitaran} $\bar B$.
  \item Tunjukkanlah $r(\bar B) \subseteq S$ dan $r(x) = x$ bagi $x
        \in S$ \emph{(untuk $\norm x = 1$, periksalah $t(x) = 0$
        dengan memakai $\langle x, u(x)\rangle \geq 0$, yang sendirinya
        menyusul dari $\langle x, x - g(x)\rangle = 1 -
        \langle x, g(x)\rangle \geq 0$)}. Simpulkanlah bersama
        Bagian II: \emph{setiap pemetaan diri mulus $\bar B$ mempunyai
        titik tetap.}
\end{enumerate}

\medskip
\noindent\textbf{Bagian IV --- Brouwer \omterm{def:b3:topology:continuity}{kontinu}.} Misalkan
$f\colon\bar B\to\bar B$ \omterm{def:b3:topology:continuity}{kontinu} tanpa titik tetap.
\begin{enumerate}[resume]
  \item Tunjukkanlah $\varepsilon = \min_{\bar B}\norm{f - x} > 0$,
        lalu hasilkanlah pemetaan \emph{polinomial} $p\colon\R^n\to
        \R^n$ dengan $\sup_{\bar B}\norm{p - f} <
        \varepsilon/2$ (Stone--Weierstrass,
        \cref{thm:b3:complete:stoneweierstrass}, koordinat
        demi koordinat --- benarkanlah peralihan dari hampiran
        skalar ke hampiran vektor).
  \item Pemetaan $p$ boleh jadi meninggalkan bolanya; tetapkanlah $g = \frac{p}{1
        + \varepsilon/2}$. Tunjukkanlah $g(\bar B) \subseteq \bar B$
        dan $\sup_{\bar B}\norm{g - f} < \varepsilon$.
  \item Turunkanlah kontradiksi dengan Bagian III lalu simpulkanlah:
        \emph{setiap \omterm{def:b3:topology:continuity}{pemetaan kontinu} $\bar B^n \to \bar B^n$
        mempunyai titik tetap.}
  \item Tunjukkanlah lewat contoh bahwa teoremanya gagal pada: bola
        terbuka; sfera $S$; dan anulus \omterm{def:b3:forms:closedexact}{tertutup}. Sifat
        $\bar B$ manakah yang hilang pada masing-masing contoh tandingnya?
\end{enumerate}

\medskip
\noindent\textbf{Bagian V --- Dividennya.}
\begin{enumerate}[resume]
  \item (Perron--Frobenius, keberadaan) Misalkan $A$ matriks
        $n\times n$ yang semua entrinya $> 0$, dan
        $\Delta = \{x \in \R^n : x_i \geq 0,\ \sum x_i =
        1\}$. Tunjukkanlah bahwa pemetaan $x \mapsto Ax/\norm{Ax}_1$
        terdefinisi baik dan \omterm{def:b3:topology:continuity}{kontinu} pada $\Delta$, bahwa
        $\Delta$ \omterm{def:b3:topology:continuity}{homeomorfik} dengan bola \omterm{def:b3:forms:closedexact}{tertutup}
        $\R^{n-1}$ (lewat \omterm{def:b3:topology:continuity}{homeomorfisma} radial dari \omterm{def:b3:topology:compact}{kompak}
        cembung yang \omterm{def:b3:topology:interior}{interiornya} tak kosong pada rentang afinnya),
        lalu simpulkanlah bahwa $A$ mempunyai vektor eigen
        berentri positif sejati dan nilai eigen $> 0$.
  \item Simpulkanlah bahwa setiap matriks stokastik berentri
        positif (yang kolomnya berjumlah $1$) mempunyai vektor
        peluang stasioner $\pi = A\pi$ --- yakni
        vektor bertipe PageRank. (Ketunggalannya juga berlaku tetapi
        memerlukan alat lain.)
  \item (Bola berbulu, persiapan) Misalkan $v$ mulus pada
        \omterm{def:b3:topology:topology}{persekitaran} $S = S^{n-1}$ dengan $\langle v(x),
        x\rangle = 0$ dan $\norm{v(x)} = 1$ bagi $x \in S$
        (yakni \emph{medan singgung satuan}). Untuk $t \in \R$ tetapkan
        $F_t(x) = x + t\,v(x)$. Tunjukkanlah
        $\norm{F_t(x)} = \sqrt{1 + t^2}$ pada $S$: sehingga $F_t$ memetakan
        $S$ ke dalam sfera $\sqrt{1+t^2}\,S$.
  \item Tunjukkanlah bahwa $P(t) = \int_SF_t^*\sigma$ merupakan
        \emph{polinomial} dalam $t$ \emph{(sebab setiap fungsi
        koefisien $F_t^*\sigma$ pada sebuah peta bersifat polinomial
        dalam $t$, dengan koefisien yang mulus terhadap variabel
        petanya; sedangkan pengintegralannya linear)}.
  \item Tunjukkanlah bahwa untuk $\abs t$ yang kecil, $F_t$ merupakan
        difeomorfisma dari $S$ ke $\sqrt{1+t^2}\,S$:
        keinjektifannya untuk $t\operatorname{Lip}(v) < 1$; difeomorfisma
        lokalnya menurut teorema fungsi balikan
        (\cref{thm:b3:submanifolds:ift} pada peta); sedangkan petanya
        terbuka sekaligus \omterm{def:b3:forms:closedexact}{tertutup} pada sfera sasarannya yang \omterm{def:b3:topology:connected}{terhubung}.
        Simpulkanlah, dengan memakai \cref{lem:b3:forms:consistency} dan
        penskalaan $\sigma_{\lambda x} =
        \lambda^{n}\,\sigma_x$ di bawah $x \mapsto \lambda x$
        (periksalah), bahwa
        \[
        P(t) = \pm(1 + t^2)^{n/2}\int_S\sigma,
        \qquad\text{dengan tanda } + \text{ untuk nilai kecil } t
        \]
        (\omterm{def:b3:forms:orientation}{orientasinya} terawetkan lewat \omterm{def:b3:topology:continuity}{kekontinuan} dari $t = 0$).
  \item Simpulkanlah (Milnor): jika $n$ \emph{ganjil}, maka $(1 +
        t^2)^{n/2}$ bukan polinomial dalam $t$, padahal ia
        berimpit dengan polinomial $P(t)/\int_S\sigma$ di dekat
        $0$ --- yakni kontradiksi. Karena itu \emph{sfera
        berdimensi genap $S^{n-1}$ (dengan $n$ ganjil) tak mengemban
        medan singgung satuan}, dan, lewat penormalan dan
        pemulusan (konvolusikan komponen demi komponen lalu proyeksikan ---
        benarkanlah kedua langkahnya), tak ada pula medan singgung \omterm{def:b3:topology:continuity}{kontinu}
        yang tak lenyap di mana pun: sehingga setiap angin di Bumi menyisakan
        satu titik tenang.

  \item (Sfera ganjil tersisir bebas) Tunjukkanlah pada $S^{2m-1}
        \subseteq \R^{2m} \cong \C^m$ sebuah medan singgung
        satuan mulus yang eksplisit: $v(x) = \iu x$ dalam notasi
        kompleksnya, yakni
        \[
        v(x_1, y_1, \dots, x_m, y_m) =
        (-y_1, x_1, \dots, -y_m, x_m) .
        \]
        Periksalah kesinggungannya dan panjang satuannya, lalu simpulkanlah bahwa
        dikotomi paritas pertanyaan 23 bersifat tajam: bahwa sebuah
        sfera tersisir tepat saat dimensinya
        ganjil. Di manakah argumen polinomial pertanyaan
        22 patah untuk $n$ yang genap?
  \item (Nol dari perilaku batasnya) Misalkan $f \colon \bar
        B^n \to \R^n$ \omterm{def:b3:topology:continuity}{kontinu} dengan $\langle f(x),
        x\rangle \geq 0$ bagi setiap $x \in S^{n-1}$. Tunjukkanlah
        bahwa $f$ lenyap di suatu tempat di $\bar B^n$.
        \emph{(Sebab jika tidak, $g(x) = -f(x)/\norm{f(x)}$ memetakan
        $\bar B$ secara \omterm{def:b3:topology:continuity}{kontinu} ke $S \subseteq \bar B$;
        lalu terapkanlah Brouwer pada $g$ dan bertentanganlah dengan hipotesis
        batasnya.)} Simpulkanlah \emph{kriteria
        kesurjektifannya}: bahwa $F\colon\R^n\to\R^n$ yang \omterm{def:b3:topology:continuity}{kontinu} dengan
        $\frac{\langle F(x), x\rangle}{\norm x} \to
        +\infty$ saat $\norm x \to \infty$ bersifat surjektif ---
        yakni leluhur berdimensi berhingga bagi argumen koersivitas
        pada analisis taklinear.

\end{enumerate}
\end{problem}
